\documentclass[11pt,a4paper]{article}
\usepackage[margin=1in]{geometry}
\usepackage{hyperref}
\usepackage{xcolor}
\usepackage{listings}
\usepackage{booktabs}
\usepackage{enumitem}
\usepackage{fancyhdr}
\usepackage{graphicx}
\usepackage{amsmath}
\usepackage{tabularx}
\usepackage{float}

\hypersetup{colorlinks=true,linkcolor=blue,urlcolor=blue}

\lstset{
  basicstyle=\ttfamily\small,
  backgroundcolor=\color{gray!10},
  frame=single,
  framerule=0.5pt,
  rulecolor=\color{gray!50},
  breaklines=true,
  breakatwhitespace=true,
  columns=fullflexible,
  keepspaces=true,
  showstringspaces=false,
  xleftmargin=4pt,
  xrightmargin=4pt,
  aboveskip=8pt,
  belowskip=8pt,
}

\pagestyle{fancy}
\fancyhf{}
\fancyhead[L]{\texttt{sheLLaMa}}
\fancyhead[R]{Cloud AI Cost Estimation}
\fancyfoot[C]{\thepage}

\title{\textbf{sheLLaMa as a Cloud AI Cost Estimation Platform}\\[4pt]\large Measuring Local LLM Usage to Project Cloud Provider Spend}
\author{sheLLaMa Project}
\date{May 2026}

\begin{document}
\maketitle
\thispagestyle{fancy}

\begin{abstract}
Organizations adopting AI face a critical question: what would our workload cost on cloud AI providers? sheLLaMa answers this by running inference locally on commodity hardware while simultaneously tracking every token consumed and projecting costs across 28+ cloud models from OpenRouter, Amazon Bedrock, OpenAI, Anthropic, Google, xAI, and Meta. This paper describes the architecture, methodology, and practical use of sheLLaMa as a cloud AI cost estimation and capacity planning tool.
\end{abstract}

\tableofcontents
\newpage

%----------------------------------------------------------------------
\section{Introduction}

Cloud AI services charge per token---typically split into input (prompt) and output (completion) tokens. Pricing varies by orders of magnitude: from \$0.035/M tokens (Amazon Nova Micro) to \$75/M tokens (Claude Opus 4 output). For organizations evaluating AI adoption, the question is not whether AI is useful, but whether the cost is justified at their usage volume.

sheLLaMa is an AI-powered agentic shell that integrates directly into bash and PowerShell. The AI reads files, writes code, runs commands, searches the web, and iterates---with permission tiers, conversation memory, and context awareness. It supports local LLMs (Ollama), cloud fallback, and distributed backends. Crucially, it also doubles as a cloud AI cost estimation platform:

\begin{enumerate}
  \item Run your actual AI workloads locally (code generation, chat, analysis, shell automation)
  \item The platform counts every prompt and response token consumed
  \item Real-time cost projections show what the same usage would cost on each cloud provider
  \item Per-client and per-API-key cost attribution enables team-level budgeting
  \item Benchmarking compares local model quality against cloud alternatives
\end{enumerate}

The result is a data-driven cost estimate based on \emph{actual usage patterns}, not hypothetical projections.

%----------------------------------------------------------------------
\section{The Cost Estimation Problem}

\subsection{Why Estimation is Hard}

Cloud AI pricing is complex:

\begin{itemize}
  \item \textbf{Asymmetric pricing}: Input and output tokens have different rates (e.g., Claude 4 Sonnet: \$3/M input, \$15/M output)
  \item \textbf{Variable output length}: The same prompt can produce vastly different response lengths depending on the model and task
  \item \textbf{Workload mix}: Organizations use AI for diverse tasks---chat, code generation, document analysis, shell automation---each with different token profiles
  \item \textbf{Rapid price changes}: Cloud providers update pricing frequently; static spreadsheet estimates become stale
  \item \textbf{Hidden costs}: Rate limiting, retries, conversation context accumulation, and cache misses all inflate real-world spend
\end{itemize}

\subsection{sheLLaMa's Approach}

Rather than estimating from assumptions, sheLLaMa measures actual usage:

\begin{itemize}
  \item Every request records \texttt{prompt\_tokens} and \texttt{response\_tokens} from Ollama's response metadata
  \item Token counts are accumulated per client IP, per task type, and per API key
  \item Cost projections use \textbf{live pricing} fetched from OpenRouter's API and AWS Pricing API
  \item Static fallback prices are maintained for when APIs are unreachable
  \item Benchmark tokens are excluded from cost projections to reflect real usage only
\end{itemize}

%----------------------------------------------------------------------
\section{Architecture for Cost Tracking}

\subsection{Token Recording Pipeline}

Every request flowing through the frontend load balancer is instrumented:

\begin{lstlisting}[language=Python,caption={Token recording in proxy\_request}]
record_ip_tokens(
    ip=client_ip,
    tokens=result.get('total_tokens', 0),
    task_type=task_type,
    prompt_tokens=result.get('prompt_tokens', 0),
    response_tokens=result.get('response_tokens', 0),
    cloud_fallback=result.get('cloud_fallback', False),
    key_name=get_key_name()
)
\end{lstlisting}

Tokens are stored in three dimensions:

\begin{itemize}
  \item \textbf{Cumulative totals} (persisted to disk, survive restarts)
  \item \textbf{Per-IP history} (timestamped entries for time-range queries)
  \item \textbf{Per-task breakdown} (chat, generate, explain, analyze, image)
\end{itemize}

\subsection{Pricing Data Sources}

\begin{table}[H]
\centering
\begin{tabularx}{\textwidth}{lXl}
\toprule
\textbf{Source} & \textbf{Models Covered} & \textbf{Refresh} \\
\midrule
OpenRouter API & Claude, GPT-4o, Gemini, Grok, Llama, Nova (15 models) & Per /test request \\
AWS Pricing API & Bedrock on-demand: Nova, Llama, DeepSeek, Mistral (9 models) & Per /test request \\
Static fallback & All 28 models (hardcoded) & Code updates \\
\bottomrule
\end{tabularx}
\caption{Pricing data sources and coverage}
\end{table}

The system attempts live pricing first. If OpenRouter or AWS is unreachable (offline environments), static fallback prices are used. The response always includes \texttt{pricing\_source} (\texttt{openrouter} or \texttt{static}) so users know the data freshness.

\subsection{Cost Calculation}

For each provider, the cost formula is:

\[
\text{Cost} = \frac{\text{prompt\_tokens} \times P_{\text{input}}}{1{,}000{,}000} + \frac{\text{response\_tokens} \times P_{\text{output}}}{1{,}000{,}000}
\]

where $P_{\text{input}}$ and $P_{\text{output}}$ are per-million-token rates.

%----------------------------------------------------------------------
\section{Supported Cloud Providers and Models}

\subsection{OpenRouter Models (via API)}

\begin{table}[H]
\centering
\begin{tabular}{lrr}
\toprule
\textbf{Model} & \textbf{Input (\$/M)} & \textbf{Output (\$/M)} \\
\midrule
Claude 4 Sonnet & 3.00 & 15.00 \\
Claude 4 Haiku & 1.00 & 5.00 \\
Claude 3.5 Sonnet & 3.00 & 15.00 \\
GPT-4o & 2.50 & 10.00 \\
GPT-4o mini & 0.15 & 0.60 \\
OpenAI o3 & 2.00 & 8.00 \\
OpenAI o4-mini & 1.10 & 4.40 \\
Gemini 2.5 Pro & 1.25 & 10.00 \\
Gemini 2.5 Flash & 0.30 & 2.50 \\
Grok 3 & 3.00 & 15.00 \\
Grok 3 mini & 0.30 & 0.50 \\
Llama 3.1 70B & 0.40 & 0.40 \\
Amazon Nova Pro & 0.80 & 3.20 \\
Amazon Nova Lite & 0.06 & 0.24 \\
Amazon Nova Micro & 0.04 & 0.14 \\
\bottomrule
\end{tabular}
\caption{OpenRouter model pricing (static fallback values; live prices may differ)}
\end{table}

\subsection{Amazon Bedrock Models (via AWS Pricing API)}

\begin{table}[H]
\centering
\begin{tabular}{lrr}
\toprule
\textbf{Model} & \textbf{Input (\$/M)} & \textbf{Output (\$/M)} \\
\midrule
Claude Opus 4 & 15.00 & 75.00 \\
Claude 4 Sonnet & 3.00 & 15.00 \\
Claude 3.5 Sonnet & 3.00 & 15.00 \\
Claude 3.5 Haiku & 0.80 & 4.00 \\
Nova Pro & 0.80 & 3.20 \\
Nova Lite & 0.06 & 0.24 \\
Nova Micro & 0.035 & 0.14 \\
Nova Premier & 2.50 & 12.50 \\
Llama 4 Maverick & 0.24 & 0.97 \\
Llama 4 Scout & 0.17 & 0.66 \\
Llama 3.3 70B & 0.72 & 0.72 \\
DeepSeek R1 & 1.35 & 5.40 \\
Mistral Large 3 & 0.50 & 1.50 \\
\bottomrule
\end{tabular}
\caption{Amazon Bedrock on-demand pricing (us-east-1, static fallback values)}
\end{table}

\subsection{Enterprise Subscriptions (Seat-Based)}

Unlike per-token APIs, enterprise AI services charge a flat monthly fee per user. sheLLaMa compares these by calculating what your actual token usage would cost on the underlying per-token API, then showing the difference against the subscription price.

\begin{table}[H]
\centering
\begin{tabular}{llr}
\toprule
\textbf{Service} & \textbf{Underlying Model} & \textbf{\$/seat/month} \\
\midrule
Microsoft Copilot & GPT-4o & 30.00 \\
ChatGPT Team & GPT-4o, o3 & 30.00 \\
ChatGPT Enterprise & GPT-4o, o3 (unlimited) & 60.00 \\
Claude Team & Claude 4 Sonnet, Opus 4 & 30.00 \\
Claude Enterprise & Claude 4 Sonnet (higher limits) & 60.00 \\
Google Gemini Business & Gemini 2.5 Pro & 24.00 \\
Google Gemini Enterprise & Gemini 2.5 Pro (unlimited) & 36.00 \\
Grok (X Premium+) & Grok 3 & 40.00 \\
Amazon Q Developer Pro & Bedrock Claude 4 Sonnet & 19.00 \\
Amazon Q Business Pro & Bedrock Nova Pro & 20.00 \\
\bottomrule
\end{tabular}
\caption{Enterprise AI subscription pricing (monthly per seat)}
\end{table}

The comparison answers: ``Given our team's actual token consumption, is a per-seat subscription cheaper or more expensive than pay-per-token API access?''

For light users, subscriptions are expensive relative to actual API consumption. For heavy users who would exceed API budgets, subscriptions provide cost certainty. The \texttt{/enterprise-costs?users=N\&days=N} endpoint calculates this dynamically.

%----------------------------------------------------------------------
\section{Cost Estimation Interfaces}

\subsection{Admin Console --- Costs Page}

The web-based costs page (\texttt{/costs}) provides:

\begin{itemize}
  \item \textbf{Time range filtering}: Day, week, month, year, all-time, or custom date range
  \item \textbf{Hypothetical costs}: What your total local usage would cost on each cloud provider
  \item \textbf{Actual fallback spend}: Real costs incurred when cloud fallback was triggered
  \item \textbf{Cached savings}: Tokens served from cache that would have cost money on cloud
  \item \textbf{Visual indicators}: Cheapest provider highlighted green, most expensive in red
  \item \textbf{Bedrock section}: Separated with AWS orange styling for clarity
  \item \textbf{Live pricing source}: Shows whether prices are from OpenRouter or static fallback
\end{itemize}

\subsection{Admin Console --- Cloud Usage Page}

The cloud usage page (\texttt{/cloud-usage}) provides per-client cost attribution:

\begin{itemize}
  \item \textbf{By Client IP}: Token usage and projected cloud cost per client, with expandable provider breakdown
  \item \textbf{By API Key}: Same breakdown per authenticated API key
  \item \textbf{By Cloud Model}: Which fallback models were used and their token consumption
  \item \textbf{Cost range}: Each row shows cheapest and most expensive provider estimate
  \item \textbf{Expandable detail}: Click to see full provider-by-provider cost table per client
\end{itemize}

The underlying API endpoint (\texttt{/cloud-usage-data}) returns JSON with cost estimates computed using the same pricing pipeline as the costs page.

\subsection{CLI Benchmarking}

The \texttt{,test} command provides per-request cost estimates:

\begin{lstlisting}[caption={CLI benchmark with cost estimates}]
$ ,test all
Testing 4 models with default prompt...

Model               Time    Prompt  Response  Total   tok/s
qwen2.5-coder:3b    12.3s   142     387       529     43.1
qwen2.5-coder:7b    34.7s   142     412       554     16.0
deepseek-coder:6.7b 38.2s   142     395       537     14.1
qwen2.5-coder:14b   97.1s   142     445       587     6.0

Cloud cost estimates (avg 142 prompt + 410 response tokens):
Provider            Input       Output      Total
Claude 4 Sonnet     $0.000426   $0.006150   $0.006576
GPT-4o              $0.000355   $0.004100   $0.004455
Amazon Nova Micro   $0.000005   $0.000057   $0.000062
...
\end{lstlisting}

\subsection{REST API}

\begin{lstlisting}[caption={Cost estimation via API}]
# Get cumulative cost projections
curl http://server:5000/cloud-costs

# Get costs for a specific time range
curl "http://server:5000/cost-history?since=1715900000&until=1716000000"

# Benchmark with cost comparison
curl -X POST http://server:5000/test \
  -H "Content-Type: application/json" \
  -d '{"model": "all"}'
\end{lstlisting}

%----------------------------------------------------------------------
\section{Use Cases for Cost Estimation}

\subsection{Pre-Migration Planning}

Before committing to a cloud AI provider, run sheLLaMa for 1--4 weeks with your team's actual workload. The costs page will show:

\begin{itemize}
  \item Total monthly projected spend per provider
  \item Which provider offers the best price for your workload mix
  \item Whether cheaper models (Nova Micro, GPT-4o mini) are viable alternatives
  \item The ratio of input to output tokens in your usage (affects provider choice)
\end{itemize}

\subsection{Budget Forecasting}

The time-range filtering enables trend analysis:

\begin{itemize}
  \item Daily token consumption patterns (peak hours, quiet periods)
  \item Week-over-week growth as AI adoption increases
  \item Per-user or per-team breakdown via API key tracking
  \item Seasonal patterns in development cycles
\end{itemize}

\subsection{Provider Comparison}

The 28-model comparison table answers questions like:

\begin{itemize}
  \item ``Is Bedrock cheaper than OpenRouter for our workload?''
  \item ``How much would we save switching from Claude to Nova Pro?''
  \item ``What's the cost difference between GPT-4o and GPT-4o mini for our use case?''
  \item ``Is the 5x price premium of Claude Opus justified over Claude Sonnet?''
\end{itemize}

\subsection{ROI Calculation for Local Infrastructure}

Compare the projected cloud spend against the cost of local hardware:

\begin{table}[H]
\centering
\begin{tabularx}{\textwidth}{Xrr}
\toprule
\textbf{Scenario} & \textbf{Local Cost} & \textbf{Cloud Cost (monthly)} \\
\midrule
Light usage (10K tokens/day) & Server: \$2,000 one-time & \$0.50--\$15/mo \\
Medium usage (100K tokens/day) & Server: \$4,000 one-time & \$5--\$150/mo \\
Heavy usage (1M tokens/day) & Server: \$8,000 one-time & \$50--\$1,500/mo \\
Team (10 devs, 5M tokens/day) & 2 servers: \$10,000 & \$250--\$7,500/mo \\
\bottomrule
\end{tabularx}
\caption{Break-even analysis: local hardware vs.\ cloud AI spend}
\end{table}

At heavy usage, local infrastructure pays for itself in 1--6 months depending on the cloud provider chosen.

\subsection{Hybrid Strategy Optimization}

sheLLaMa's cloud fallback feature enables a hybrid approach:

\begin{enumerate}
  \item Run most requests locally (zero marginal cost)
  \item Fall back to cloud only when local quality is insufficient
  \item Track actual fallback spend separately from hypothetical costs
  \item Set per-key budgets to cap cloud spend: \texttt{"budget": \{"max\_daily": 5.00\}}
  \item Receive webhook alerts at 80\% budget consumption
\end{enumerate}

The costs page separates hypothetical costs (``what if everything ran on cloud'') from actual fallback spend (``what we actually paid''), giving a clear picture of savings.

%----------------------------------------------------------------------
\section{Implementation Details}

\subsection{Persistent Statistics}

Cost data survives service restarts:

\begin{itemize}
  \item Cumulative token totals saved to \texttt{persisted\_totals.json} every 60 seconds
  \item Per-IP history maintained in memory with configurable retention
  \item Background thread handles periodic persistence
  \item \texttt{/reset-cloud-costs} endpoint for admin to clear data (requires admin role)
\end{itemize}

\subsection{Benchmark Exclusion}

Tokens from \texttt{/test} benchmark requests are excluded from cost projections. This ensures the running tab reflects real usage only---not artificial load from model comparison tests.

\subsection{Cache Awareness}

When prompt caching is enabled (default 5-minute TTL), cached responses are tracked separately. The costs page shows how many tokens were served from cache---representing additional savings that would not exist on most cloud providers (which charge per request regardless of caching).

\subsection{Live Pricing Refresh}

\begin{lstlisting}[language=Python,caption={Live pricing fetch from OpenRouter}]
def fetch_cloud_pricing():
    resp = requests.get('https://openrouter.ai/api/v1/models', timeout=10)
    data = resp.json().get('data', [])
    by_id = {}
    for m in data:
        p = m.get('pricing', {})
        prompt = float(p.get('prompt', '0'))
        comp = float(p.get('completion', '0'))
        if prompt > 0 or comp > 0:
            by_id[m['id']] = {'input': prompt * 1e6, 'output': comp * 1e6}
    # Map display names to live prices...
\end{lstlisting}

Bedrock pricing is fetched via the AWS CLI:

\begin{lstlisting}[caption={AWS Pricing API call}]
aws pricing get-products --service-code AmazonBedrock \
  --region us-east-1 \
  --filters Type=TERM_MATCH,Field=regionCode,Value=us-east-1
\end{lstlisting}

\subsection{Per-Key Budget Enforcement}

API keys can have spending limits that are enforced on actual cloud fallback spend:

\begin{lstlisting}[language={},caption={API key with budget limit}]
{
  "api_keys": {
    "sk-team-dev": {
      "name": "dev-team",
      "role": "user",
      "budget": {"max_daily": 10.00},
      "rate_limit": {"rpm": 30, "tpd": 500000}
    }
  }
}
\end{lstlisting}

When the budget is 80\% consumed, a webhook fires. At 100\%, cloud fallback requests are rejected with HTTP 429.

%----------------------------------------------------------------------
\section{Deployment for Cost Estimation}

\subsection{Minimal Setup (Single Server)}

\begin{lstlisting}[language=bash,caption={Quick start for cost estimation}]
# Install and start
ollama pull qwen2.5-coder:7b
pip install flask ollama pyyaml requests psutil
python backend/app.py

# Point your team's CLI clients at the server
export SHELLAMA_API=http://server:5000
export SHELLAMA_MODEL=qwen2.5-coder:7b
source cli/shellama.bash

# Use AI normally for 1-4 weeks, then check costs
curl http://server:5000/cloud-costs | python -m json.tool
\end{lstlisting}

\subsection{Distributed Setup (Production)}

For larger teams, deploy the frontend load balancer with multiple backends:

\begin{lstlisting}[language={},caption={backends.json for multi-server deployment}]
{
  "backends": [
    {"url": "http://192.168.1.230:5000", "weight": 10, "max_model": "qwen2.5-coder:7b"},
    {"url": "http://192.168.1.233:5000", "weight": 10, "max_model": "qwen2.5-coder:14b"}
  ]
}
\end{lstlisting}

All token tracking is centralized on the frontend, regardless of which backend processes the request.

%----------------------------------------------------------------------
\section{Interpreting Results}

\subsection{Reading the Cost Table}

A typical \texttt{/cloud-costs} response after one week of team usage:

\begin{lstlisting}[language={}]
{
  "prompt_tokens": 847293,
  "response_tokens": 2341567,
  "total_tokens": 3188860,
  "pricing_source": "openrouter",
  "cloud_costs": [
    {"provider": "Claude 4 Sonnet", "total_cost": 37.66},
    {"provider": "GPT-4o", "total_cost": 25.53},
    {"provider": "Gemini 2.5 Flash", "total_cost": 6.11},
    {"provider": "Amazon Nova Micro", "total_cost": 0.36},
    {"provider": "Bedrock Claude Opus 4", "total_cost": 188.22}
  ]
}
\end{lstlisting}

\textbf{Interpretation}: This team's weekly AI usage would cost between \$0.36 (Nova Micro) and \$188.22 (Claude Opus 4) on cloud providers. The local infrastructure cost is \$0 marginal after hardware purchase.

\subsection{Token Ratio Analysis}

The ratio of prompt to response tokens affects provider choice:

\begin{itemize}
  \item \textbf{High output ratio} (code generation, explanations): Providers with low output pricing win (Nova Micro, Llama, GPT-4o mini)
  \item \textbf{High input ratio} (analysis of large files, context-heavy chat): Providers with low input pricing win (Gemini with large context, Nova Lite)
  \item \textbf{Balanced}: Overall total cost is the deciding factor
\end{itemize}

\subsection{Quality vs.\ Cost Tradeoff}

The \texttt{,test} benchmark lets you compare local model quality against cloud alternatives. If a \$0.04/request local model produces acceptable output for 95\% of tasks, the ROI of local infrastructure is clear. The 5\% that needs cloud fallback can be budgeted precisely using the actual fallback tracking.

\subsection{Subscription vs.\ Pay-Per-Token}

The enterprise comparison reveals whether seat-based subscriptions make financial sense:

\begin{lstlisting}[language=bash,caption={Enterprise cost comparison API}]
# Compare for a 10-person team over 30 days of tracked usage
curl "http://server:5000/enterprise-costs?users=10&days=30"
\end{lstlisting}

\begin{lstlisting}[language={}]
{
  "enterprise_costs": [
    {"provider": "Microsoft Copilot", "monthly_total": 300.00,
     "equivalent_token_cost": 2.25, "savings_vs_subscription": 297.75},
    {"provider": "Amazon Q Developer Pro", "monthly_total": 190.00,
     "equivalent_token_cost": 1.80, "savings_vs_subscription": 188.20}
  ]
}
\end{lstlisting}

\textbf{Interpretation}: If the subscription costs \$300/month but the equivalent API usage is only \$2.25, the team is paying 133x more for the subscription than they would on per-token pricing. This is common for light-to-moderate usage. Heavy users (millions of tokens/day) may find subscriptions cheaper due to unlimited usage caps.

%----------------------------------------------------------------------
\section{Limitations}

\begin{itemize}
  \item \textbf{Token counts are model-dependent}: Different tokenizers produce different counts for the same text. Local Ollama token counts are approximate proxies for cloud provider counts.
  \item \textbf{Pricing changes}: Live pricing is fetched per-request for benchmarks but cached otherwise. Prices may change between fetches.
  \item \textbf{No latency costing}: Cloud providers may offer faster responses; the cost of developer wait time is not captured.
  \item \textbf{No quality weighting}: A cheap model that requires 3 retries may cost more than an expensive model that succeeds first try. This is partially addressed by the fallback tracking.
  \item \textbf{Region-specific}: Bedrock pricing is for us-east-1. Other regions may differ.
\end{itemize}

%----------------------------------------------------------------------
\section{Conclusion}

sheLLaMa transforms local LLM infrastructure into a cloud AI cost estimation platform. By running actual workloads locally and projecting costs across 28+ cloud models with live pricing, organizations can make data-driven decisions about:

\begin{itemize}
  \item Whether to adopt cloud AI services
  \item Which provider offers the best value for their specific workload
  \item When local infrastructure ROI exceeds cloud spend
  \item How to budget for hybrid local+cloud deployments
  \item Per-team and per-project cost allocation
\end{itemize}

The platform requires no cloud accounts, no API keys, and no ongoing costs to operate as a cost estimator. It runs fully offline after the initial model pull, making it suitable for air-gapped environments and organizations evaluating AI adoption without committing to vendor contracts.

%----------------------------------------------------------------------
\section{References}

\begin{itemize}
  \item OpenRouter pricing API: \url{https://openrouter.ai/api/v1/models}
  \item AWS Bedrock pricing: \url{https://aws.amazon.com/bedrock/pricing/}
  \item Ollama: \url{https://ollama.com}
  \item sheLLaMa source: \url{https://github.com/rory/shellama}
\end{itemize}

\end{document}
